\documentclass[11pt,a4paper]{article}

\usepackage[utf8]{inputenc}
\usepackage[T1]{fontenc}
\usepackage[english]{babel}
\usepackage[a4paper,margin=2.4cm]{geometry}
\usepackage{amsmath,amssymb}
\usepackage{bm}
\usepackage{booktabs}
\usepackage{enumitem}
\usepackage{xcolor}
\usepackage{hyperref}
\usepackage{tcolorbox}

\hypersetup{colorlinks=true, linkcolor=blue!50!black, urlcolor=blue!60!black}

\newcommand{\BHF}{B_{\mathrm{hf}}}
\newcommand{\dEQ}{\Delta E_Q}
\newcommand{\nuR}{\nu}
\newcommand{\code}[1]{\texttt{\small #1}}

\newtcolorbox{nota}[1][Note]{colback=blue!4, colframe=blue!45!black,
  fonttitle=\bfseries, title={#1}, left=4pt, right=4pt, top=3pt, bottom=3pt}
\newtcolorbox{aviso}[1][Warning]{colback=orange!7, colframe=orange!70!black,
  fonttitle=\bfseries, title={#1}, left=4pt, right=4pt, top=3pt, bottom=3pt}

\title{\textbf{Magnetic Relaxation Models in Fitbauer}\\[3pt]
\large Phenomenological component, Blume--Tjon model, and N\'{e}el--Arrhenius size distribution}
\author{Department of Physical Chemistry · UAM}
\date{\today}

\begin{document}
\maketitle

\begin{abstract}\noindent
This document describes the three magnetic relaxation models implemented
in Fitbauer, complementary to the pre-existing singlet, doublet, sextet, and
Hesse--R\"{u}bartsch distribution models:
(1)~\code{Relajacion}: phenomenological component that interpolates between a
blocked sextet and an apparent superparamagnetic doublet;
(2)~\code{BlumeTjon}: dynamic two-state model ($+\BHF \leftrightarrow -\BHF$)
with symmetric Bloch--McConnell-type exchange;
(3)~\code{NeelSize}: N\'{e}el--Arrhenius model with a lognormal particle-size
distribution, combining the \code{BlumeTjon} profile for each diameter.
Simultaneous global fitting of spectra at multiple temperatures is available
via the \code{GlobalNeelFitDialog} (\code{gui/dialogs.py},
\textbf{Fit $\to$ Global N\'{e}el (T)}).
\end{abstract}

\tableofcontents
\bigskip

%======================================================================
\section{Physical motivation}

In magnetic nanoparticles, ferrites, or partially superparamagnetic iron
oxides, the magnetic moment may fluctuate within the M\"{o}ssbauer observation
window. If the fluctuations are slow, the nucleus observes a nearly static
hyperfine field and the spectrum is a sextet. If they are fast, the magnetic
field is averaged out and a superparamagnetic doublet or singlet appears. In the
intermediate regime the lines broaden and coalesce.

The implementation aims to avoid automatic interpretation of the type
``sextet + doublet = two chemical phases''. A doublet-like fraction may be the
spectral manifestation of fast relaxation.

\begin{aviso}
The models described here are complementary. If the \code{Relajacion} or
\code{BlumeTjon} components are not selected, the earlier models are evaluated
through exactly the same code branches as before.
\end{aviso}

%======================================================================
\section{Common notation}

Transmission is written, in the thin-absorber limit, as
\begin{equation}
  T(v)=b+s\,v-A(v),
\end{equation}
where $b$ is the baseline, $s$ the slope, and $A(v)$ the positive absorption of
the component. For a conventional sextet:
\begin{equation}
  A_{6}(v)=D\sum_{j=1}^{6} I_j L\bigl(v; v_j,\Gamma_j\bigr),
\end{equation}
with depth $D$, intensities $I_j$, and Lorentzian profile
\begin{equation}
  L(v;v_0,\Gamma)=\frac{\Gamma^2}{(v-v_0)^2+\Gamma^2}.
\end{equation}
The first-order positions are
\begin{equation}
  v_j = \delta + q_j\,\dEQ + m_j\,\frac{\BHF}{33\,\mathrm{T}},
\end{equation}
where $m_j$ are the $\alpha$-Fe reference positions at $33$ T and $q_j$ the
first-order quadrupolar pattern.

%======================================================================
\section{Phenomenological model \code{Relajacion}}

\subsection{Definition}

The \code{Relajacion} component interpolates between a blocked sextet and an
apparent superparamagnetic doublet:
\begin{equation}
  A_{\mathrm{rel}}(v)=D\left[
  f_{\mathrm{b,ef}}\,S_6(v)
  +\bigl(1-f_{\mathrm{b,ef}}\bigr)\,D_2^{\ast}(v)
  \right].
  \label{eq:rel-phen}
\end{equation}
$S_6(v)$ is the unit-depth sextet profile and $D_2^{\ast}(v)$ is the doublet
scaled so that its integrated area is comparable to that of the sextet:
\begin{equation}
  D_2^{\ast}(v)=D_2(v)\,
  \frac{\int S_6(v)\,dv}{\int D_2(v)\,dv}.
\end{equation}
Thus $D$ approximately retains the meaning of total area/depth.

\subsection{Blocked fraction and phenomenological rate}

The user sees two specific parameters:
\begin{itemize}[leftmargin=*]
  \item $f_{\mathrm{bloq}}$: maximum blocked fraction.
  \item $\log_{10}\nu$: phenomenological relaxation rate in s$^{-1}$.
\end{itemize}
The fraction that actually enters \eqref{eq:rel-phen} is
\begin{equation}
  f_{\mathrm{b,ef}}=f_{\mathrm{bloq}}\,F(\log_{10}\nu),
  \label{eq:fb-eff}
\end{equation}
with
\begin{equation}
  F(x)=\frac{1}{1+\exp\left(\frac{x-x_c}{w}\right)},
  \qquad x_c=8.5,\quad w=0.55.
\end{equation}
Hence slow rates ($x\ll x_c$) leave $F\simeq 1$, while fast rates ($x\gg x_c$)
force $F\simeq 0$.

In addition, a maximum phenomenological broadening is applied in the
intermediate regime:
\begin{equation}
  \Gamma_{\mathrm{ef}}=\Gamma\left[1+1.6\exp\left(-\frac{1}{2}
  \left(\frac{x-x_c}{0.75}\right)^2\right)\right].
\end{equation}

\begin{aviso}[Important correlation]
$f_{\mathrm{bloq}}$ and $\log_{10}\nu$ are not fully independent parameters:
both reduce the observable sextet fraction. In practice, it is advisable to fix
one of the two unless the spectrum contains sufficient information. To estimate
fractions, fix $\log_{10}\nu$ and free $f_{\mathrm{bloq}}$; to study dynamics,
fix $f_{\mathrm{bloq}}=1$ and free $\log_{10}\nu$.
\end{aviso}

%======================================================================
\section{Dynamic model \code{BlumeTjon}: two states}

\subsection{Exchange matrix}

The \code{BlumeTjon} component implements a first dynamic version with two
opposite field orientations:
\begin{equation}
  +\BHF \leftrightarrow -\BHF,
\end{equation}
with symmetric transition matrix
\begin{equation}
  W=\begin{pmatrix}-\nu & \nu\\ \nu & -\nu\end{pmatrix}.
\end{equation}
There is no explicit blocked fraction: the dynamics are controlled by $\nu$.

\subsection{Exchange profile for a single transition}

For each line of the sextet two frequencies/velocities are computed, one for
$+\BHF$ and one for $-\BHF$, denoted $v_a$ and $v_b$. A symmetric
Bloch--McConnell/Kubo exchange lineshape is used. We define
\begin{equation}
  z_a = \Gamma + i(v-v_a),\qquad z_b = \Gamma + i(v-v_b),
\end{equation}
and an exchange rate in velocity units
\begin{equation}
  k_v = \frac{\nu}{\mathcal{C}_{\nu\to v}},
\end{equation}
where
\begin{equation}
  \mathcal{C}_{\nu\to v}=\frac{E_\gamma}{h c_{\mathrm{mm/s}}}.
\end{equation}
Here $E_\gamma$ is the energy of the $^{57}$Fe transition and
$c_{\mathrm{mm/s}}$ the speed of light in mm/s. The absorption profile used
is the real part of
\begin{equation}
  R(v)=\frac{1}{2}\,\Gamma\,
  \frac{z_a+z_b+4k_v}{z_a z_b+k_v(z_a+z_b)}.
  \label{eq:exchange}
\end{equation}
The total contribution is
\begin{equation}
  A_{\mathrm{BT}}(v)=D\sum_{j=1}^{6} I_j\,\Re\left[R_j(v)\right]_+,
\end{equation}
where $[\cdot]_+$ clips small numerical negative contributions.

\subsection{Limits}

\begin{itemize}[leftmargin=*]
  \item \textbf{Slow limit}: $k_v\to0$. Equation \eqref{eq:exchange}
  tends to the average of two static Lorentzians. In the symmetric
  $+\BHF/-\BHF$ pattern it reproduces the blocked sextet.
  \item \textbf{Intermediate}: $k_v$ comparable to the line separation.
  Broadening and coalescence appear.
  \item \textbf{Fast}: large $k_v$. The magnetic positions are averaged and
  the spectrum tends to a doublet or singlet depending on $\dEQ$.
\end{itemize}

%======================================================================
\section{N\'{e}el--Arrhenius and size distribution}

The \code{NeelSize} component uses the N\'{e}el--Arrhenius model to convert
particle volume to relaxation rate. For spherical particles of diameter $d$,
\begin{equation}
  V(d)=\frac{\pi d^3}{6},
\end{equation}
and
\begin{equation}
  \tau_N(d,T)=\tau_0\exp\left(\frac{K_{\mathrm{eff}}V(d)}{k_B T}\right),
  \qquad
  \nu(d,T)=\tau_0^{-1}\exp\left(-\frac{K_{\mathrm{eff}}V(d)}{k_B T}\right).
\end{equation}
In the GUI, $\log_{10}K_{\mathrm{eff}}$ and $\log_{10}\tau_0$ are fitted to
maintain reasonable numerical ranges:
\begin{equation}
  \log_{10}\nu_i=-\log_{10}\tau_0
  -\frac{K_{\mathrm{eff}}V(d_i)}{k_BT\ln 10}.
\end{equation}

The size distribution is initially modelled as lognormal. If $d_{50}$ is the
median and $\sigma_d$ the lognormal width,
\begin{equation}
  f(d)=\frac{1}{d\sigma_d\sqrt{2\pi}}
  \exp\left[-\frac{(\ln d-\ln d_{50})^2}{2\sigma_d^2}\right].
\end{equation}
The total spectrum is discretised as
\begin{equation}
  A_{\mathrm{Neel}}(v)=D\sum_i w_i\,
  A_{\mathrm{BT}}\bigl(v;\log_{10}\nu_i,T,K_{\mathrm{eff}},d_i\bigr),
\end{equation}
where $A_{\mathrm{BT}}$ is the two-state exchange profile from the preceding
section. Large particles have a higher barrier $K_{\mathrm{eff}}V$ and therefore
relax more slowly; small particles relax fast and contribute superparamagnetic
signal.

The approximate blocking temperature for a given diameter is estimated as
\begin{equation}
  T_B(d)=\frac{K_{\mathrm{eff}}V(d)}{k_B\ln(\tau_M/\tau_0)},
\end{equation}
where $\tau_M$ is the M\"{o}ssbauer observation time (by default $10^{-8}$~s
is used for the GUI summary).

\begin{aviso}[Identifiability]
In a single spectrum, $K_{\mathrm{eff}}$, $d_{50}$, $\sigma_d$, $\tau_0$, $\Gamma$,
and a possible static distribution $P(\BHF)$ can be strongly correlated. It is
therefore advisable to fix parameters using external information
(TEM, magnetometry, literature) or to use multi-temperature global fitting.
\end{aviso}

%======================================================================
\section{Multi-temperature global fitting}

Phase 5 adds a pure backend for simultaneously fitting several spectra at
different temperatures. For each spectrum $r$, there is $T_r(v)$ and temperature
$T_r$. The model shares the physical parameters
\begin{equation}
  \theta_{\mathrm{shared}}=\{\delta,\dEQ,\BHF,\Gamma,K_{\mathrm{eff}},d_{50},
  \sigma_d,\tau_0\},
\end{equation}
and maintains local per-spectrum parameters, by default
\begin{equation}
  \theta_r=\{b_r,s_r,D_r\}.
\end{equation}
The objective function is
\begin{equation}
  \chi^2=\sum_r\sum_k
  \left(\frac{T_{r,\mathrm{model}}(v_k;\theta_{\mathrm{shared}},\theta_r,T_r)
  -T_{r,\mathrm{data}}(v_k)}{\sigma_{rk}}\right)^2.
\end{equation}
This reduces the degeneracy compared with fitting a single spectrum: the same
size distribution must simultaneously explain the blocking at low $T$ and the
superparamagnetic collapse at high $T$.

%======================================================================
\section{Relation to fitting and reports}

The adjustable parameters are the same as in a sextet ($\delta$, $\dEQ$,
$\BHF$, linewidths, depth, and intensities) plus $\log_{10}\nu$ and, only in the
phenomenological model, $f_{\mathrm{bloq}}$. In \code{NeelSize}, $T$,
$\log_{10}K_{\mathrm{eff}}$, $d_{50}$, $\sigma_d$, $\log_{10}\tau_0$, and the
number of bins are added. The GUI also shows
\begin{equation}
  \tau=\frac{1}{\nu}=10^{-\log_{10}\nu}\;\mathrm{s}
\end{equation}
and, for \code{NeelSize}, an approximate $T_B$.

In reports, the \code{Relajacion} component should be interpreted as
``blocked / apparent superparamagnetic fraction'', not as separate chemical
phases. The \code{BlumeTjon} component is reported as a dynamic two-state model
and \code{NeelSize} as a N\'{e}el--Arrhenius model with a size distribution.

%======================================================================
\section{Implemented validation}

Synthetic tests were added to verify:
\begin{enumerate}[leftmargin=*]
  \item \code{Relajacion} with $f_{\mathrm{bloq}}=1$ reproduces a sextet.
  \item \code{Relajacion} with $f_{\mathrm{bloq}}=0$ reproduces a doublet
  scaled to the sextet area.
  \item $\log_{10}\nu$ controls the slow/fast limits of the phenomenological
  model.
  \item \code{BlumeTjon} reproduces the slow limit and fast exchange
  qualitatively.
  \item The lognormal size distribution is correctly normalised and the N\'{e}el
  rate decreases as diameter increases.
  \item \code{NeelSize} changes with temperature and the global backend
  recovers a shared synthetic diameter from a pair of spectra.
\end{enumerate}

\begin{nota}[Scope — v4.9.0+]
The three models described in this document (\code{Relajacion}, \code{BlumeTjon},
and \code{NeelSize}) are fully implemented in \code{core/physics.py} and
integrated in the GUI. Multi-temperature global fitting is available from
\textbf{Fit $\to$ Global N\'{e}el (T)} (\code{GlobalNeelFitDialog} in
\code{gui/dialogs.py}). For a single spectrum it is advisable to fix parameters
using external information (TEM, magnetometry) or to use multi-temperature global
fitting to reduce degeneracy.
\end{nota}

\end{document}
